%% kapa: exam=1 round=35 year=2024 period=1 form=A subject=real_estate subject_ko=부동산학원론 seq=13 number=93
%% source: Q-Net 감정평가사 기출문제 (cst00309) artlSeq=5213229 file=raw/1008/5213229/1교시 문제지 원본.pdf
%% conversion: kordoc 4.24.51 -> tools/gen1.py (auto); review=proofread-sonnet; flags=none
\begin{question}{93}
원가방식에 관한 설명으로 옳은 것을 모두 고른 것은?
\begin{framedtext}
ㄱ. 원가법과 적산법은 원가방식에 해당한다. \par
ㄴ. 재조달원가는 실제로 생산 또는 건설된 방법 여하에 불구하고 도급방식을 기준으로 산정한다. \par
ㄷ. 대상부동산이 가지는 물리적 특성인 지리적 위치의 고정성에 의해서 경제적 감가요인이 발생한다. \par
ㄹ. 정액법, 정률법, 상환기금법은 대상부동산의 내용연수를 기준으로 하는 감가수정방법에 해당한다. \par
\end{framedtext}
\begin{choices}
  \item ㄱ, ㄴ
  \item ㄷ, ㄹ
  \item ㄱ, ㄴ, ㄹ
  \item ㄱ, ㄷ, ㄹ
  \item ㄱ, ㄴ, ㄷ, ㄹ
\end{choices}
\end{question}
\answer{5}
